\documentclass[11pt]{article}
\usepackage[margin=1.15in]{geometry}
\usepackage{amsmath,amssymb,amsthm}
\usepackage{xcolor}
\definecolor{linkblue}{rgb}{0.1,0.15,0.5}
\usepackage[colorlinks=true,linkcolor=linkblue,citecolor=linkblue,urlcolor=linkblue]{hyperref}

\newtheorem{theorem}{Theorem}
\newtheorem{proposition}{Proposition}
\newtheorem{corollary}{Corollary}
\newcommand{\Reff}{R_{\mathrm{eff}}}
\newcommand{\E}{\mathbb{E}}
\newcommand{\Var}{\mathrm{Var}}

\title{Rating Gaps and Effective Resistance:\\
Exact Fluctuation Laws for Elo and Bradley--Terry Systems}
\author{The \textit{fixed-gaps} authors\thanks{Code, data pipelines, and the
companion monograph are available at
\url{https://github.com/macrokit/fixed-gaps}. This paper develops Chapter 7
of the monograph.}}
\date{July 2026}

\begin{document}
\maketitle

\begin{abstract}
A rating system is a network: players are nodes, matches are edges, and
ratings are node potentials whose differences predict outcomes. We show
that the fluctuation behavior of rating \emph{gaps} is governed by a single
question --- \emph{where noise enters the network relative to where matches
defend it} --- and that the answer splits rating systems into three regimes
with three different laws. (i)~\emph{Estimation}: the variance of the
maximum-likelihood estimate of any rating gap equals the effective
resistance between the two players in the match graph, with Fisher
conductances. (ii)~\emph{Dynamics with static skills}: the stationary gap
fluctuation of online Elo is $K$ (the update factor) for \emph{every} pair
on \emph{every} connected graph --- outcome noise and restoring force enter
through the same edges and their graph dependence cancels exactly.
(iii)~\emph{Dynamics with drifting skills}: idiosyncratic skill drift
restores the resistance law as the tracking error, with mean reversion
acting as a grounding conductance on every node. All three laws are
verified in simulation; the estimation law is additionally confirmed on
up to 104{,}419 real Chatbot Arena battles with a zero-parameter fit
(slope $1.035$--$1.041$ against the theoretical $1$, $R^2 \ge 0.958$), in
both dense and clustered battle graphs. The clustered experiment --- two
collection eras joined by five bridge models --- confirms the framework's
distinctive prediction: the resistance law stays calibrated where the
battle-count heuristic drifts by $17\%$, and quantifies how the
comparability of $103$ models across a year of LLM history rests on five
survivors. Applications include resistance-based leaderboard confidence
intervals, a matchmaking design rule, and a structural explanation of
chess's era problem.
\end{abstract}

\section{Introduction}

Pairwise-comparison rating systems --- Elo in chess \cite{elo1978}, its
Bradley--Terry foundation \cite{zermelo1929,bradleyterry1952}, Glicko
\cite{glickman1999}, TrueSkill \cite{trueskill2006}, and most recently
LLM leaderboards such as Chatbot Arena \cite{chiang2024} --- share one
structure: a set of entities, a graph of realized comparisons, and a
vector of ratings whose pairwise \emph{differences} are the only
predictive content. The absolute level of a rating is pure convention (the
familiar incomparability of ratings across pools); the gaps are the
physics.

This paper studies the fluctuations of those gaps and shows they are
organized by a single structural question:
\begin{quote}
\emph{Where does randomness enter the comparison network, relative to
where the matches act?}
\end{quote}
Three answers give three laws. Randomness in \emph{outcomes} enters on the
edges --- the same edges through which every match pulls the rating gap
toward the observed result. Randomness in \emph{true skills} (players
improving and declining) enters at the nodes. And randomness \emph{common}
to all players enters nowhere: it cancels from every gap identically.

\paragraph{Contributions.}
\begin{enumerate}
\item \textbf{Estimation resistance law} (Thm.~\ref{thm:est}). The
variance of the maximum-likelihood estimate of any rating gap --- between
any two players, whether or not they ever met --- equals the effective
resistance between them in the match graph with Fisher conductances. The
ingredients are classical; the identification of the inverse Fisher
information as a resistance, and the consequences drawn from it, appear
not to have been exploited in the rating literature.
\item \textbf{Isotropy of online Elo} (Thm.~\ref{thm:iso}). With static
skills, the stationary fluctuation of the rating gap under online Elo
with update factor $K$ is $\Var(r_i - r_j) = K$ for every pair, on every
connected match graph: a path and a complete graph produce identical
long-run churn. To our knowledge this graph-independence is unobserved in
the literature; it is invisible unless one asks the noise-location
question, since it results from an exact cancellation between edge noise
and edge defense.
\item \textbf{Tracking resistance law} (Thm.~\ref{thm:track}). When true
skills drift idiosyncratically (variance rate $\sigma_d^2$ per node), the
resistance law returns as the stationary \emph{tracking} error:
$\Var\!\big((r_i - r_j) - (s_i - s_j)\big) = K + \tfrac{2|E|\sigma_d^2}{K}
\Reff(i,j)$, and mean reversion of skills at rate $\theta$ replaces
$\Reff$ by the resistance of the network in which every node is grounded
through conductance $\theta/a$. Without mean reversion, the logistic
saturates and no bounded-update system can track at all
(Prop.~\ref{prop:sat}).
\item \textbf{Validation.} All laws verified in controlled simulation
(\S\ref{sec:sim}); the estimation law verified on real Chatbot Arena data
with no free parameters, in a dense regime (39{,}241 battles, 60 models,
slope $1.035$, $R^2=0.958$) and in a clustered two-era regime (104{,}419
battles, 108 models, slope $1.041$, $R^2 = 0.963$), where the framework's
distinctive prediction --- calibration of the resistance prediction
where the count heuristic fails --- is confirmed (\S\ref{sec:arena}).
\item \textbf{Design consequences} (\S\ref{sec:design}): resistance-based
confidence intervals for leaderboards; a matchmaking rule (bridge
comparisons dominate: one cross-cluster match outperformed a
within-cluster match by a factor ${\sim}44$ on worst-pair resistance in
our test configuration); and a derivation of chess's era-comparison
problem as a chain instability, together with an explanation of why
Arena's era problem is mild.
\end{enumerate}

\section{Setup}\label{sec:setup}

Players $1,\dots,n$ carry true skills $s \in \mathbb{R}^n$; player $i$
beats $j$ with probability $\sigma(s_i - s_j)$, $\sigma(x) = (1 +
e^{-x})^{-1}$ (Bradley--Terry; Elo's base-$10$/$400$ scale is a linear
change of units). The \emph{match graph} $G = (V, E)$ records which pairs
play; $L$ denotes its Laplacian and, for conductances $c_e$ on the edges,
$L_c$ the weighted Laplacian. The \emph{effective resistance} between $i$
and $j$ is $\Reff(i,j) = (e_i - e_j)^\top L_c^{+} (e_i - e_j)$ with
$L_c^{+}$ the pseudoinverse \cite{doylesnell1984}. Online Elo updates,
after a match on edge $(i,j)$ with score $x \in \{0,1\}$:
\[
r_i \mathrel{+}= K\big(x - \sigma(r_i - r_j)\big), \qquad
r_j \mathrel{-}= K\big(x - \sigma(r_i - r_j)\big),
\]
which is stochastic gradient ascent on the Bradley--Terry log-likelihood
\cite{mandt2017}. Matches arrive uniformly on $E$; linearizing about
small gaps ($\sigma'(0) = \tfrac14$) gives the drift--diffusion
\begin{equation}\label{eq:lin}
dr = a\,L\,(s - r)\,dt + dM, \qquad a = \frac{K}{4|E|}, \qquad
\mathrm{Cov}(dM) = \frac{K^2}{4|E|}\,L\,dt .
\end{equation}
The structural fact driving everything below: \emph{the noise covariance
and the restoring force are built from the same Laplacian}, because in a
rating system every observation of a gap is also a defense of it.

\section{Three laws}\label{sec:laws}

\begin{theorem}[Estimation resistance law]\label{thm:est}
With $n_e$ matches on edge $e$ and static skills, the Fisher information
of the Bradley--Terry model is $L_c$ with $c_e = n_e\, p_e(1-p_e)$, $p_e =
\sigma(s_i - s_j)$. Hence, asymptotically, for every pair $(i,j)$,
\[
\Var\big(\widehat{r_i - r_j}\big) \;=\; \Reff(i,j)
\quad\text{in the match graph with Fisher conductances;}
\]
at equal skills and $m$ matches per edge this is $(4/m)\,
\Reff^{\mathrm{unit}}(i,j)$.
\end{theorem}

\begin{proof}
The log-likelihood Hessian at the truth is $-\sum_e n_e\, p_e(1-p_e)(e_i -
e_j)(e_i - e_j)^\top = -L_c$; the MLE covariance on the mean-zero gauge is
$L_c^{+}$, and the variance of the contrast $(e_i - e_j)^\top \hat r$ is
the stated quadratic form, i.e.\ the resistance \cite{ghoshboydsaberi2008}.
\end{proof}

Uncertainty about a gap is thus a property of the \emph{network between}
the pair --- parallel comparison paths add like parallel conductances ---
not of the two players' match counts.

\begin{theorem}[Isotropy of stationary Elo]\label{thm:iso}
Under \eqref{eq:lin} with static equal skills, the stationary covariance
of $r$ on the mean-zero subspace is $\Sigma = \tfrac{K}{2} P$ (with $P$
the centering projector). Consequently
\[
\Var(r_i - r_j) \;=\; K \qquad \text{for every pair, on every connected
match graph.}
\]
\end{theorem}

\begin{proof}
The stationary Lyapunov equation for $dr = -aLr\,dt + dM$ is $aL\Sigma +
\Sigma L a = \frac{K^2}{4|E|} L$. Since drift and noise are simultaneously
diagonalizable (both are $L$), $\Sigma = \frac{K^2/(4|E|)}{2a} P =
\frac{K}{2}P$ on the range of $L$, and $(e_i - e_j)^\top P (e_i - e_j) =
2$.
\end{proof}

The graph dependence cancels exactly: a poorly connected pair receives
less restoring force but, in the same proportion, less noise. Online
Elo's long-run churn is graph-independent --- contrary to the natural
intuition that remote pairs should wander more, and, to our knowledge,
unremarked in the rating literature.

\begin{theorem}[Tracking resistance law]\label{thm:track}
Let true skills drift idiosyncratically, $ds = \sigma_d\, dW$ (independent
across players), and let $e = r - s$ be the tracking error under
\eqref{eq:lin}. Then in stationarity
\[
\Var(e_i - e_j) \;=\; \underbrace{K}_{\text{edge noise
(Thm.~\ref{thm:iso})}} \;+\; \underbrace{\frac{2|E|\,\sigma_d^2}{K}\,
\Reff(i,j)}_{\text{node noise: the resistance law}} .
\]
If instead skills are mean-reverting, $ds = -\theta s\,dt + \sigma_d\,dW$,
the exact linear theory replaces $\Reff$ by the \emph{grounded resistance}
--- the quadratic form of $(L + (\theta/a) I)^{-1}$: mean reversion is a
grounding wire of conductance $\theta/a$ on every node.
\end{theorem}

\begin{proof}
Mode-by-mode on the eigenbasis of $L$ (eigenvalue $\lambda$): the error
obeys $de_\lambda = -a\lambda\, e_\lambda\,dt + (\text{edge noise, rate }
\tfrac{K^2}{4|E|}\lambda) - (\text{node noise, rate } \sigma_d^2)$, giving
stationary variance $\frac{K}{2} + \frac{\sigma_d^2}{2a\lambda}$; summing
$\lambda^{-1}$ against the pair contrast yields $\Reff$. With mean
reversion the transfer function from skill noise to error acquires the
factor $-i\omega/\big((i\omega + \theta)(i\omega + a\lambda)\big)$, whose
integrated power replaces $(a\lambda)^{-1}$ by $(\theta +
a\lambda)^{-1}$, i.e.\ grounds the network.
\end{proof}

\begin{proposition}[Saturation: the era catastrophe]\label{prop:sat}
If skills random-walk without mean reversion, skill gaps grow without
bound, $\sigma'(s_i - s_j) \to 0$, and the linearization gain in
\eqref{eq:lin} collapses: the tracking error of a bounded-update system
diverges on every graph.
\end{proposition}

Common shocks --- randomness applied to all players alike --- cancel from
every gap identically and are invisible to all three laws: only
\emph{asymmetric} randomness moves gaps. The trichotomy
common/node/edge is exhaustive and is the organizing principle of this
paper.

\section{Simulation}\label{sec:sim}

We verify the laws against the \emph{actual} (nonlinear, discrete) Elo
dynamics, not the linearization. Configurations: $n = 12$--$16$ players on
path, cycle, two-clique-plus-two-bridges, random 3-regular, and complete
graphs; $K = 0.1$. Code and committed output accompany the paper.

\textbf{Thm.~\ref{thm:est}}: regressing empirical MLE gap-error variance
(600 trials, 40 matches/edge) on $(4/m)\Reff$ over all pairs gives slope
$1.006$--$1.027$ (theory $1$) with $R^2$ up to $0.994$ across the graph
families (the complete graph has all resistances nearly equal, hence no
variance to explain).

\textbf{Thm.~\ref{thm:iso}}: online Elo, $4 \times 10^6$ steps, static
skills: mean stationary gap variance $0.1055$--$0.1060$ against the
predicted $K = 0.1$ (the small excess is the expected nonlinear
correction), and flat in resistance --- on the path, the low-resistance
half of pairs shows $0.1056$ and the high-resistance half $0.1064$.

\textbf{Thm.~\ref{thm:track}}: drifting skills ($\sigma_d = 0.004$,
$\theta = 10^{-4}$, $8\times10^6$ steps): tracking-error variance is
linear in $\Reff$ (path $R^2 = 0.95$, two-cluster $0.90$), intercept
$0.106/0.105$ vs.\ the predicted $K = 0.1$; measured slopes $0.00206$
(path) and $0.00964$ (two-cluster) against exact grounded-theory values
$0.00213$ (3\%) and $0.00812$ (within 20\%), with the naive ungrounded
value ($0.00352$ on the path) clearly rejected.

\textbf{Prop.~\ref{prop:sat}}: with $\theta = 0$ the path-graph tracking
variance reaches ${\sim}60\times$ the defended level and keeps growing.

\section{Real data: Chatbot Arena}\label{sec:arena}

\subsection{Dense regime}
From the public 55k Arena preference dataset we retain decisive battles
among models with $\ge 200$ of them: $39{,}241$ battles, $60$ models,
$1{,}145$ of $1{,}770$ pairs realized (density $0.65$). We fit the
Bradley--Terry MLE, compute the predicted variance of every pairwise gap
as the Fisher resistance (Thm.~\ref{thm:est}), and compare against $200$
bootstrap resamples of the battles. The theory has no free parameters:
the predicted slope is exactly $1$. Over all $1{,}770$ pairs:
\[
\text{slope} = 1.035, \qquad \text{corr} = 0.979, \qquad
R^2_{\text{through origin}} = 0.958 .
\]
The law is quantitatively right at the few-percent level, simultaneously
across every pair, including the majority that never met. On this dense
graph the folk heuristic $1/n_i + 1/n_j$ also correlates well (log--log
$0.956$ vs.\ $0.987$ for resistance): in dense graphs the resistance
approximately \emph{reduces} to the count formula --- the theory explains
when the heuristic works and predicts where it fails, namely clustered
graphs. Which motivates:

\subsection{Clustered regime: the era experiment}
Merging the 55k collection (2023--early 2024) with the 100k collection
(June--August 2024) yields a two-era graph: $104{,}419$ decisive battles,
$108$ models --- $55$ exclusive to era~1, $48$ exclusive to era~2, and a
\textbf{five-model bridge} (three GPT-4 variants, GPT-3.5-turbo, Mixtral)
through which all $55 \times 48 = 2{,}640$ cross-era comparisons must
route, since era-exclusive models never coexisted. Findings:
\begin{enumerate}
\item The resistance law holds on the clustered graph: slope $1.041$,
corr $0.982$, $R^2 = 0.963$ over all $5{,}778$ pairs.
\item Calibrating the count heuristic on within-era pairs and evaluating
the median ratio of bootstrap variance to prediction: the resistance
prediction reads $1.01$ within-era and $1.01$ cross-era ---
identically calibrated in both regimes with no adjustment --- while the
count heuristic reads $0.92$ within vs.\ $1.08$ cross, a systematic
$17\%$ spread in the predicted direction.
\item The bridge is load-bearing: deleting a single bridge model's
battles raises median cross-era resistance by up to $8\%$; deleting all
five disconnects the eras --- cross-era gaps become infinitely
uncertain. The comparability of $103$ models across a year of LLM
history rests, in this precise sense, on five survivors.
\item Arena's cross-era penalty is nonetheless mild (${\sim}8\%$ in
variance at matched counts), because its bridge models are battle-rich:
the platform keeps old flagships in rotation, which is exactly the
long-range-edge prescription of \S\ref{sec:design}. Chess, whose
bridges (individual careers spanning eras) are thin and unreplayable,
is the same graph with weak bridges --- hence its severe era problem.
\end{enumerate}

\section{Design consequences}\label{sec:design}

\paragraph{Resistance-based confidence intervals.} By
Thm.~\ref{thm:est}, leaderboard uncertainty for ``$A$ vs.\ $B$'' should
be reported as the resistance between $A$ and $B$ in the battle graph.
Two models with thousands of battles each but no common opponents can
have arbitrarily uncertain relative rank; conversely a lightly played
model wired into a well-connected neighborhood can be sharply placed.

\paragraph{Matchmaking.} To sharpen a leaderboard fastest, schedule the
comparison that most reduces the (grounded) resistance of the pairs of
interest. In our two-cluster test, one bridge match reduced worst-pair
resistance from $1.000$ to $0.778$; one more within-cluster match, to
$0.995$ --- a ${\sim}44\times$ difference per match. Platforms that
matchmake for entertainment systematically overspend on low-resistance
edges.

\paragraph{The era problem.} The temporal comparison graph of a
long-lived rating pool is chain-like; by Thm.~\ref{thm:track} cross-era
gaps drift with variance linear in era separation (and by
Prop.~\ref{prop:sat}, without anchoring, tracking eventually fails
outright). Known remedies map onto the theory: anchoring to a fixed
reference is grounding; careers spanning eras are long-range edges;
Arena's bridge-rich rotation policy is the design rule implemented.

\section{Related work}\label{sec:related}

The Bradley--Terry model and its MLE date to \cite{zermelo1929,
bradleyterry1952}; Elo's system to \cite{elo1978}; measurement-error
refinements include Glicko \cite{glickman1999} and TrueSkill
\cite{trueskill2006}. Effective resistance is classical in the
random-walk literature \cite{doylesnell1984} and appears in optimal
experiment design and network analysis \cite{ghoshboydsaberi2008}; the
interpretation of SGD's stationary distribution as an
Ornstein--Uhlenbeck process follows \cite{mandt2017}. Chatbot Arena and
its Bradley--Terry leaderboard are described in \cite{chiang2024}, which
computes bootstrap confidence intervals; our Thm.~\ref{thm:est} gives
those intervals a closed structural form, and \S\ref{sec:arena} verifies
it. We are not aware of prior statements of the isotropy law
(Thm.~\ref{thm:iso}) or of the grounded-resistance tracking law
(Thm.~\ref{thm:track}) in this setting. The general framework relating
noise location, defended invariants, and effective resistance is
developed in the companion monograph (footnote, p.~1), of which this
paper is the self-contained applied chapter.

\section{Limitations and outlook}

The dynamical results are linearizations, and the simulations quantify
the nonlinear corrections at the few-percent level for realistic $K$;
Thm.~\ref{thm:iso}'s constant, in particular, is $K$ only to leading
order. The Arena validation tests the estimation law; catching the
tracking law (Thm.~\ref{thm:track}) in the wild requires a longitudinal
pool with thin bridges --- historical chess federation data is the
natural target, but game-level archives at the required scale are not
freely available. The era experiment treats the two collections as
exchangeable judge populations; drift in judging standards would appear
as a common shock and cancel from gaps, but interactions between judge
mix and model identity would not, and bounding them is future work.

\begin{thebibliography}{9}
\bibitem{zermelo1929} E.~Zermelo.
Die Berechnung der Turnier-Ergebnisse als ein Maximumproblem der
Wahrscheinlichkeitsrechnung. \emph{Math.\ Z.}, 29:436--460, 1929.
\bibitem{bradleyterry1952} R.~A.~Bradley and M.~E.~Terry.
Rank analysis of incomplete block designs: I.\ The method of paired
comparisons. \emph{Biometrika}, 39:324--345, 1952.
\bibitem{elo1978} A.~E.~Elo.
\emph{The Rating of Chessplayers, Past and Present}. Arco, 1978.
\bibitem{glickman1999} M.~E.~Glickman.
Parameter estimation in large dynamic paired comparison experiments.
\emph{J.\ R.\ Stat.\ Soc.\ C}, 48:377--394, 1999.
\bibitem{trueskill2006} R.~Herbrich, T.~Minka, and T.~Graepel.
TrueSkill: a Bayesian skill rating system. \emph{NeurIPS}, 2006.
\bibitem{doylesnell1984} P.~G.~Doyle and J.~L.~Snell.
\emph{Random Walks and Electric Networks}. MAA, 1984.
\bibitem{ghoshboydsaberi2008} A.~Ghosh, S.~Boyd, and A.~Saberi.
Minimizing effective resistance of a graph. \emph{SIAM Review},
50(1):37--66, 2008.
\bibitem{mandt2017} S.~Mandt, M.~D.~Hoffman, and D.~M.~Blei.
Stochastic gradient descent as approximate Bayesian inference.
\emph{JMLR}, 18(134):1--35, 2017.
\bibitem{chiang2024} W.-L.~Chiang et al.
Chatbot Arena: an open platform for evaluating LLMs by human preference.
\emph{ICML}, 2024.
\end{thebibliography}

\end{document}
